\section*{Capítulo \ref{ch:b2:lab-techniques-2} --- Técnicas de laboratório II: síntese e análise na prática}

\begin{solution}{exo:b2:lab-techniques-2:1}
Gelo e água; gelo e sal; gelo-seco e acetona (como para os \omterm{def:b2:enolates-aldol:enolate}{enolatos} do \cref{ch:b2:enolates-aldol}).
\end{solution}

\begin{solution}{exo:b2:lab-techniques-2:2}
A camada inferior. Adicione algumas gotas de água: elas se juntam à camada aquosa, aqui a superior.
\end{solution}

\begin{solution}{exo:b2:lab-techniques-2:3}
As primeiras porções se aglomeram ao absorver água; quando o sólido recém-adicionado permanece solto e
gira livremente, a solução está seca.
\end{solution}

\begin{solution}{exo:b2:lab-techniques-2:4}
O $R_f$ é alto demais para uma coluna (deseja-se cerca de 0.3) e as duas manchas estão próximas: use um
eluente menos polar, com mais hexano (9\,:\,1, por exemplo), o que abaixa os dois $R_f$ e aumenta sua
separação em volumes de coluna.
\end{solution}

\begin{solution}{exo:b2:lab-techniques-2:5}
\begin{itemize}
  \item Hidrogenocarbonato de sódio (pH 8.3 $>$ 4.2 + 2): benzoato na água; acidifique, extraia o ácido benzoico.
  \item Hidróxido de sódio (pH $>$ 12 $>$ 9.99 + 2): fenóxido na água; acidifique, extraia o fenol.
  \item Ácido clorídrico (pH cerca de 1 $<$ 4.6 $-$ 2): anilínio na água; basifique, extraia a anilina.
  \item O naftaleno fica no éter: seque, evapore.
\end{itemize}
A lavagem com hidrogenocarbonato deve vir primeiro: o hidróxido extrairia juntos o ácido e o fenol.
% ledger: ph:NaHCO3, pka:benzoic, pka:phenol, pka:anilinium
\end{solution}

\begin{solution}{exo:b2:lab-techniques-2:6}
$1/T = 1/T_b - (R/\Delta_{\mathrm{vap}}H)\ln(p/p^\circ)$: bromobenzeno cerca de \qty{49}{\degreeCelsius},
benzaldeído cerca de \qty{68}{\degreeCelsius} (figura da \cref{prop:b2:lab-techniques-2:reduced-pressure}).
\end{solution}

\begin{solution}{exo:b2:lab-techniques-2:7}
$0.050 \times 200\,000/400 = \qty{25}{K}$.
\end{solution}

\begin{solution}{exo:b2:lab-techniques-2:8}
Cada composto tem sua própria resposta: a porcentagem de área só é uma porcentagem em massa se os \omterm{def:b2:chromatography:quantitative}{fatores
de resposta} forem iguais. RMN: $0.12/3 = 0.040$~mol de impureza por mol de produto, isto é, $0.040 \times 100 = \qty{4}{g}$
por \qty{200}{g} de produto; $w = 200/204 \approx 98.0~\%$.
\end{solution}

\begin{solution}{exo:b2:lab-techniques-2:9}
Massa teórica $0.0300 \times 164 = \qty{4.92}{g}$. Sem correção $4.10/4.92 \approx 83.3~\%$; com correção
$0.950 \times 4.10/4.92 \approx 79.2~\%$.
\end{solution}

\begin{solution}{exo:b2:lab-techniques-2:10}
Benzeno significa que o reagente se formou e foi depois protonado: água na vidraria, no solvente ou no
aldeído; ar (umidade) entrando por um vazamento; um aldeído parcialmente oxidado a ácido benzoico, cujo
hidrogênio ácido destrói o reagente. Remédios: vidraria seca em estufa, solvente anidro, aldeído recém-destilado, uma sobrepressão de nitrogênio verificada no borbulhador. Quase nenhum reagente indicaria
uma falha de iniciação (superfície do magnésio coberta de óxido), evitada com aparas novas ou ativadas.
\end{solution}

\begin{solution}{exo:b2:lab-techniques-2:11}
\qty{270}{kJ} a \qty{150}{W}: no mínimo $270\,000/150 = \qty{1800}{s}$, meia hora, e só se o haleto
reagir tão depressa quanto é adicionado. Mais depressa, a liberação de calor excede o resfriamento: a
temperatura sobe, ou o reagente se acumula e pode depois reagir de uma só vez
(\cref{prop:b2:lab-techniques-2:accumulation}).
\end{solution}

\begin{solution}{exo:b2:lab-techniques-2:12}
Recristalização: $13.0 \times 0.80 = \qty{10.4}{g}$ a 99~\%. Coluna: $13.0 \times 0.95 =
\qty{12.4}{g}$ a 99.5~\%, mas \qty{1.5}{L} de solvente a comprar, evaporar e descartar. A coluna dá
cerca de \qty{2}{g} a mais de um produto um pouco mais puro; a recristalização é mais simples, mais
barata e gera menos resíduos, e é preferida quando sua pureza basta.
\end{solution}

\begin{solution}{pb:b2:lab-techniques-2:1}
\textbf{1.} Etoxietano: extremamente inflamável, nocivo, forma peróxidos. Bromobenzeno: inflamável,
tóxico, tóxico para a vida aquática. Magnésio: sólido inflamável, libera hidrogênio com a água.
\textbf{2.} \ce{C6H5MgBr + H2O -> C6H6 + Mg(OH)Br}: cada gota de água destrói reagente; a secagem em
estufa remove o filme de água sobre o vidro.
\textbf{3.} Para manter fora a umidade e o oxigênio, que também destrói o reagente.
\textbf{4.} O fluxo de nitrogênio; ele impede o ar de entrar pela saída.
\textbf{5.} Ele não reage com o reagente, os pares não ligantes de seu oxigênio estabilizam o magnésio, e
seu baixo ponto de ebulição (\qty{307.7}{K}) permite que o refluxo limite a temperatura.
\textbf{6.} $15.7/157.01 = \qty{0.1000}{mol}$ de bromobenzeno; $2.67/24.305 = \qty{0.1099}{mol}$ de magnésio,
em excesso de 10~\%.
\textbf{7.} $0.1000 \times 270 = \qty{27.0}{kJ}$.
\textbf{8.} Para que o calor seja liberado na velocidade com que o refluxo e o banho o removem.
\textbf{9.} $0.0200 \times 270 = \qty{5.40}{kJ}$.
\textbf{10.} $5400/170 \approx \qty{32}{K}$.
\textbf{11.} Cerca de \qty{52}{\degreeCelsius}, muito acima do ponto de ebulição do etoxietano
(\qty{34.6}{\degreeCelsius}): uma ebulição violenta que pode sobrecarregar o condensador e lançar vapor
inflamável para fora do balão.
\textbf{12.} Adicione cerca de um décimo do bromobenzeno, espere a reação começar (turvação, refluxo
suave), depois adicione o resto na velocidade de um refluxo suave, com o banho de gelo à mão.
\textbf{13.} \ce{C6H5MgBr + C6H5CHO} dá o alcóxido \ce{(C6H5)2CHOMgBr}; o benzaldeído,
$9.55/106.12 = \qty{0.0900}{mol}$, é o limitante.
\textbf{14.} O álcool é duplamente benzílico: um ácido forte o transformaria em um cátion estabilizado,
levando a éteres ou a produtos de substituição; o cloreto de amônio é ácido o bastante para protonar o
alcóxido e dissolver os sais de magnésio.
\textbf{15.} A camada superior: o etoxietano é menos denso que a água.
\textbf{16.} Ela remove a maior parte da água dissolvida no éter.
\textbf{17.} Sulfato de magnésio até que fique solto, filtração, depois um evaporador rotativo sob
pressão reduzida.
\textbf{18.} O brometo de fenilmagnésio se acopla com o bromobenzeno que não reagiu; uma concentração local
alta de bromobenzeno e uma temperatura alta o favorecem, outra razão para uma adição lenta.
\textbf{19.} A bifenila ($R_f$ 0.85), a menos polar.
\textbf{20.} Cerca de $1/0.85 \approx 1.2$ volume de coluna para a bifenila, $1/0.25 = 4$ para o álcool.
\textbf{21.} O álcool contribui com 10 prótons \omterm{def:b2:huckel:aromatic}{aromáticos} por CH; o excesso $10.90 - 10.00 = 0.90$ vem da
bifenila, 10 prótons por molécula: 0.090 mol de bifenila por mol de álcool.
\textbf{22.} $w = 184.24/(184.24 + 0.090 \times 154.21) \approx 0.930$.
\textbf{23.} $12.50 \times 0.930 \approx \qty{11.62}{g}$.
\textbf{24.} $0.0900 \times 184.24 \approx \qty{16.58}{g}$.
\textbf{25.} $12.50/16.58 \approx 75.4~\%$.
\textbf{26.} $11.62/16.58 \approx \boldsymbol{70~\%}$ (70.1~\%), corrigido pela pureza.
\end{solution}
